\section*{第 \S\arabic{exercisegroup} 节习题}
\phantomsection
\addcontentsline{toc}{section}{第 \protect\S\arabic{exercisegroup} 节习题}

\begin{enumerate}
\def\labelenumi{\arabic{enumi})}
\tightlist
\item
  设 \(\mathbf{K}\) 是特征为 0 的域，\(U\) 和 \(V\) 是 \(\mathrm{K}[\mathrm{X}]\) 上的两个复合算子，且 \(W = VU = UV\) ；设 \((u_n), (v_n)(w_n)\) 是分别对应于 \(U, V, W\) 的阿佩尔多项式序列。证明：如果 \(p\) 是 \(U\) 的阶，那么
\end{enumerate}

\[
w _ {n} ^ {(p)} (\mathrm{X}) = \sum_ {k = 0} ^ {n} \binom {n} {k} v _ {n - k} (\mathrm{X}) V _ {0} (u _ {k})
\]

\[
w _ {n} (\mathrm{X} + \mathrm{Y}) = \sum_ {k = 0} ^ {n} \binom {n} {k} u _ {k} (\mathrm{X}) v _ {n - k} (\mathrm{Y}).
\]

\begin{enumerate}
\def\labelenumi{\arabic{enumi})}
\setcounter{enumi}{1}
\tightlist
\item
  设 \(\mathbf{E}\) 是 \(\mathbf{C}\) 上由函数 \(x^n\) （ \(n \in \mathbf{N}\) ）、\(e^{\lambda x}\) （ \(\lambda \in \mathbf{C}\) ，\(\lambda \neq 0\) ）与 \(|x + \mu|\) （ \(\mu \in \mathbf{R}\) ）生成的向量空间。
\end{enumerate}

\begin{enumerate}
\def\labelenumi{\alph{enumi})}
\item
  证明上述函数构成 E 的一组基。
\item
  设 \(U\) 是由下列条件在 \(\mathbf{E}\) 上定义的复合算子：\(U_0^\xi (\xi^n) = (n!)^2\) ，\(U(e^{\lambda x}) = e^{\lambda x}\) （对 \(\lambda \neq 0\) ），\(U(|x + \mu|) = |x + \mu|\) （对 \(\mu \in \mathbf{R}\) ）。指示函数 \(u(\lambda)\) 是常数 1 ，但以 \(n!\lambda^n\) 为通项的级数对任何值 \(\lambda \neq 0\) 都不收敛。
\item
  设 \(V\) 是由下列条件在 \(\mathbf{E}\) 上定义的复合算子：\(V(x^n) = x^n\) ，\(V(|x + \mu|) = |x + \mu|\) ，\(V(e^{\lambda x}) = 0\) （对 \(\lambda \neq 0\) ）；指示函数 \(v(\lambda)\) 当 \(\lambda = 0\) 时等于 1 ，当 \(\lambda \neq 0\) 时等于 0 ，因而异于级数 \(\sum_{n=0}^{\infty} \alpha_n \lambda^n\) 的和，其中 \(\alpha_n = V_0^\xi(\xi^n)/n!\) 。
\item
  设 W 是由下式在 E 上定义的复合算子
\end{enumerate}

\[
W (x ^ {n}) = x ^ {n}, \qquad W (e ^ {\lambda x}) = e ^ {\lambda x}, \qquad W (| x + \mu |) = e ^ {x + \mu};
\]

证明 \(VW \neq WV\) 。

\begin{enumerate}
\def\labelenumi{\arabic{enumi})}
\setcounter{enumi}{2}
\tightlist
\item
  设 K 是特征为 0 的域。称 K 上两个未定元 X, Y 的多项式所成的代数 \(K[X,Y]\) 的自同态 U 为复合算子，如果对每个多项式 \(f \in K[X,Y]\) ，置 \(g = U(f)\) 后有 \(g(X + S, Y + T) = U_{X,Y}(f(X + S, Y + T))\) ，其中 S 和 T 是另外两个未定元。
\end{enumerate}

\begin{enumerate}
\def\labelenumi{\alph{enumi})}
\item
  把 prop. 1 与 th. 1 推广到这些算子。由此导出公式 \(e^{\mathrm{S + T}} = e^{\mathrm{S}}e^{\mathrm{T}}\) 的一个新证明。
\item
  对形如 \(D_{X}^{p}D_{Y}^{q}u(D_{X},D_{Y})\) （其中形式级数 u 的常数项不为零）的复合算子，定义阿佩尔多项式 \(u_{mn}\) ，并推广 VI, p.~273 与 274 的 props. 4、5 和 6 。
\item
  特别考虑由 \(U(f(\mathrm{X},\mathrm{Y})) = f(\mathrm{X} + 1, \mathrm{Y} + 1) - f(\mathrm{X},\mathrm{Y} + 1) - f(\mathrm{X} + 1, \mathrm{Y}) + f(\mathrm{X},\mathrm{Y})\) 定义的复合算子 U ；对应于这个算子的阿佩尔多项式称为伯努利多项式，记作 \(B_{m,n}\) 。证明 \(\mathrm{B}_{m,n}(\mathrm{X},\mathrm{Y}) = \mathrm{B}_{m}(\mathrm{X})\mathrm{B}_{n}(\mathrm{Y})\) 。
\end{enumerate}

\setcounter{exercisegroup}{1}
\refstepcounter{exercisegroup}
